\documentclass{article}
\usepackage[margin=1in]{geometry}
\usepackage{amsmath,amssymb,amsthm}
\usepackage{booktabs}
\usepackage{graphicx}
\usepackage[numbers]{natbib}
\usepackage[colorlinks=true,linkcolor=blue,citecolor=blue]{hyperref}

\newtheorem{proposition}{Proposition}
\newcommand{\cons}{\varepsilon_{\parallel}}
\newcommand{\drift}{\varepsilon_{\perp}}

\title{Constraints Deserve No Credit:\\
Output Parameterization Governs Out-of-Distribution Error\\
in Learned Simulators with Exact Invariants}
\author{Anonymous}
\date{Draft skeleton, 2026-08-29. Families A, B, C, M adjudicated (confirmatory, frozen);
contraction mechanism test (family E) exploratory.}

\begin{document}
\maketitle

\begin{abstract}
Embedding exact invariants---conservation laws in physical surrogates, accounting identities
in financial market models---in learned simulators is a widespread practice, and reported
benefits are routinely credited to the constraint. We audit that attribution. Across three
controlled families (MLP and U-Net surrogates of advection, diffusion, and Burgers dynamics;
U-Nets under amplitude and resolution shift; and a learned simulator of a continuous double
auction with exact cash/inventory settlement), we compare six enforcement mechanisms at
matched capacity, training length, and in-distribution error, and decompose
out-of-distribution error into a conserving channel that no invariant can correct and a
constraint-violating channel. In every adjudicated setting (30/30 to date), the celebrated
out-of-distribution benefit is produced by the output parameterization (residual versus
absolute next-state prediction), not by the constraint: exact architectural enforcement leaves
conserving-channel error statistically unchanged while zeroing invariant violations, and
post-hoc projection reproduces a closed-form decoupling theorem to machine precision (42/42
settings). Strong soft penalties are dominated: they inflate conserving-channel error by
18--33\% under deterministic dynamics and by 100--370\% under stochastic market dynamics---a
domain-dependence we attribute to irreducibly stochastic targets. We distill the audit into
prescriptions: enforce invariants architecturally on residual updates (free and exact), avoid
strong penalty forms, and report the error decomposition rather than the constraint residual
alone.
\end{abstract}

\section{Introduction}
The practice and its assumption (constraint $\Rightarrow$ better simulator): climate emulators
\citep{beucler2019,beucler2021}, structure-preserving Hamiltonian networks
\citep{greydanus2019}, no-arbitrage and martingale-constrained financial models
\citep{cohen2023}. The confound: enforcing a linear invariant in a network is almost always
accompanied by a second change---predicting state \emph{increments} rather than absolute
states. No prior work separates the two at matched controls
\citep[see also][]{suh2022,pdebench}. We ask: \emph{when the invariant is enforced exactly,
what does the constraint itself contribute to out-of-distribution error, and what does the
parameterization contribute?}

\paragraph{Contributions.}
(1) A matched-pair attribution design with an exact error decomposition
(conserving vs.\ constraint-violating channels) and a built-in negative control (post-hoc
projection) that we prove and verify to be exact. (2) A cross-domain audit: three physical
families, one financial family, two architectures, four out-of-distribution probe classes.
(3) Three findings: parameterization carries the benefit; exact constraints are free but
creditless; strong soft penalties are dominated, with damage scaling in target stochasticity.

\section{Related work}
Constraint mechanisms and their evaluation in-distribution \citep{beucler2021}; OOS climate
generalization without matched parameterization controls \citep{beucler2019}; Hamiltonian
structure preservation \citep{greydanus2019}; learned PDE surrogates and OOD protocols
\citep{pdebench,pfaff2021}; gradient-access attribution in differentiable simulators
\citep{suh2022}; invariant-constrained financial models \citep{cohen2023}; laboratory market
benchmarks with exact settlement \citep{godesunder1993}. Fragment-occupancy map and the exact
gap statement are maintained in the project evidence manifest; the bounded-search novelty
statement appears here.

\section{Design}
\paragraph{Systems and invariants.} Family A: MLP one-step maps of advection, diffusion,
Burgers (periodic grid, spectral truth; invariant: mean state). Family B: two-level 1D U-Nets
at resolution 128, probed additionally at amplitude extrapolation and a band-limited $2\times$
resolution shift. Family C: 2D advection--diffusion with a 2D U-Net (pending). Family M: a
Gode--Sunder-class continuous double auction with exact settlement; invariants are the
accounting identities $\sum_i \mathrm{inv}_i = S_0$ and $\sum_i \mathrm{cash}_i = C_0$; the
learner predicts per-agent state transitions; OOD probes are high-imbalance and volatility
shock regimes.

\paragraph{Arms.} free (absolute output); free\_res (residual output); soft/soft\_res
($\lambda\in\{3,30\}$ penalty on invariant residuals); hard (residual update projected onto
the invariant tangent---exact enforcement); projection (post-hoc correction of free).

\paragraph{Decomposition and decoupling.}
For error $e = f(u) - A(u)$ and invariant projector $P$ onto the violating direction,
$\drift = Pe$ and $\cons = (I-P)e$.
\begin{proposition}[Decoupling of post-hoc projection]
For a linear invariant and its least-squares post-hoc correction, conserving-channel error is
invariant under the correction; only $\drift$ is removed.
\end{proposition}
This is the pipeline's negative control: any deviation in the projection arm indicates an
implementation fault. It held to $0.00\%$ in $42/42$ adjudicated settings.

\paragraph{Protocol.} Frozen before outcomes: capacity grids, seeds ($\ge 5$ per cell),
matched-pair same-cell comparisons, decision rules (attribution ratio $\ge 2$;
laundering conjunction; decoupling within $1\%$), budget caps, and stop rules. The
$\pm5\%$ matched-ID window triggered the pre-registered non-overlap branch (ID-cost finding)
in all families; cross-parameterization claims are therefore matched-pair only.

\section{Results}
\subsection{Family A (MLP, three PDE systems)}
Attribution ratio $260$--$5225\times$ (rule: $\ge 2$); hard $\equiv$ free\_res on
$\cons$ in all six settings; projection exact; laundering conjunction $2/6$ (inflation
without drift benefit $\Rightarrow$ dominated).

\subsection{Family B (1D U-Nets, amplitude and resolution OOD)}
Attribution $12.7$--$6710\times$ across $18/18$ settings including the $2\times$ resolution
shift; hard statistically indistinguishable from free\_res; decoupling exact $18/18$; strong
soft penalties dominated ($+22$--$30\%$ inflation, no drift benefit).

\subsection{Family M (learned CDA simulator, financial)}
Attribution $140\times$ at the matched cell; the accounting constraint's entire effect is
exact drift zeroing; \emph{laundering conjunction holds}: soft$_{30}$ inflates $\cons$ by
$4.7\times$ while buying drift reduction---damage amplifies with stochastic targets
(deterministic: $+18$--$33\%$; stochastic: $+100$--$370\%$). The absolute-output arm degrades
with capacity on noisy targets.

\subsection{Family C (2D U-Net, contractive dynamics)}
The parameterization effect \emph{inverts}: absolute output beats residual by $4.6$--$6.3\times$
on conserving-channel OOD error (e.g.\ $0.389\pm0.042$ vs $1.780\pm0.152$), and the hard
constraint tracks its residual twin ($1.862\pm0.154$). The frozen operator annihilates
high-wavenumber content in one step: for contractive maps the residual head must learn a large
cancelling update whose off-support extrapolation is fragile, while an absolute head directly
represents the smooth next state. This triggered the pre-registered freeze-and-review clause.

\subsection{Family E (mechanism review, exploratory)}
Single-factor $\nu$-sweep of the diffusion family shows a monotone dose--response: at
$\nu=0.02$ (near-identity) free\_res wins by $100\times$; at $\nu=2.0$ by $1.5\times$; at
$\nu=10.0$ (rich-band decay $e^{-4.9}$) the effect \emph{flips} and free wins by $1.9\times$
(non-overlapping seed bands at every dose; hard $\equiv$ free\_res throughout). Contraction
strength of the true one-step map is a usable predictor for choosing the parameterization.

\section{Prescriptions}
Enforce exact invariants architecturally on residual updates (free, exact); avoid strong soft
penalties; report decompositions, not residuals.

\section{Limitations and non-claims}
Linear invariants only (energy-type untested); constructed truths, no field or transfer
claims; bounded-search novelty statement; attribution statements are about error structure
under the frozen probes, not about conservation laws determining dynamics.

\section{Reproducibility}
Frozen protocols, seeds, per-run JSON records, and analyzer implementing the decision rules
verbatim are released with the code; cross-machine bit-identical smoke runs are recorded.

\bibliographystyle{plainnat}
\bibliography{refs}
\end{document}
